\title{Group Sequence Policy Optimization}

\begin{abstract}
We present Group Sequence Policy Optimization (GSPO), a reinforcement learning algorithm that provides a robust, effective, and resource-efficient approach for training large language models. Departing from earlier methods that utilize token-wise importance ratios, GSPO computes importance ratios from entire sequence probabilities and applies its clipping, reward computation, and updates at the sequence level. Our empirical results indicate that GSPO consistently outperforms the GRPO algorithm in both sample efficiency and overall model quality, offering enhanced stability during Mixture-of-Experts (MoE) RL fine-tuning and introducing the possibility of streamlining RL system architecture. The strengths of GSPO have notably driven advancements in the performance of recent Qwen3 models.
\end{abstract}

\section{Introduction}
\label{sec:intro}

Reinforcement learning (RL) has become a central methodology for enhancing the scalability of language models \citep{o1,dpsk-r1,qwq32b,qwen3}. By applying RL at scale, language models acquire the ability to address complex tasks—including advanced mathematics and programming challenges—through extended and more intricate reasoning strategies.

For RL scaling to succeed as computational resources increase, it is essential to ensure that the training process remains consistent and resilient. Nevertheless, leading RL techniques such as GRPO \citep{grpo} are plagued by significant stability concerns during the training of massive language models. Such instability frequently leads to drastic and irreversible degradation of model performance, commonly referred to as collapse \citep{qwen3, minimax-m1}, thus impeding progress toward further improving language model capabilities through prolonged RL optimization.

This study reveals that GRPO's instability arises from a fundamental error in how importance sampling weights are applied and interpreted within its algorithmic framework. This misuse leads to elevated variance in training noise, which accumulates as the model generates longer responses and is further intensified by the clipping mechanism, ultimately causing the model to fail.

To overcome these challenges, we introduce \textbf{Group Sequence Policy Optimization (GSPO)}, a novel reinforcement learning approach tailored for large language model training. GSPO's principal advancement is its rigorous and theoretically consistent importance ratio, formulated on sequence likelihood \citep{click}, thereby faithfully adhering to the principles of importance sampling. Moreover, GSPO derives normalized rewards by treating them as advantages across multiple generated responses for each prompt, establishing a strong correspondence between reward assignment at the sequence level and optimization procedure.

Our experimental results clearly show that GSPO markedly outperforms GRPO in terms of training stability, operational efficiency, and overall performance. Notably, GSPO inherently addresses stability issues commonly encountered in the RL training of large Mixture-of-Experts (MoE) models, thereby obviating the necessity for elaborate stabilization mechanisms and indicating promise for streamlining RL system design. These strengths of GSPO have directly facilitated substantial enhancements in the recent Qwen3 models. Looking forward, we anticipate GSPO serving as a resilient and scalable methodological basis that supports the further development of large-scale RL training in language modeling.

\section{Preliminaries}

\paragraph{Notation}
Throughout this work, we represent an autoregressive language model with parameters $\theta$ as a policy $\pi_\theta$. A query is symbolized by $x$, and the collection of queries is noted as $\mathcal{D}$. For any query $x$, its corresponding response $y$ has a probability assigned by the policy, written as $\pi_\theta(y | x)=\prod_{t=1}^{|y|} \pi_\theta (y_t | x, y_{<t} )$, where $|y|$ specifies the length of $y$ in tokens. The pairing $(x, y)$ is evaluated by a verifier $r$, yielding a reward value $r(x, y) \in [0, 1]$.

\paragraph{Proximal Policy Optimization (PPO)} Using samples generated from the old policy $\pi_{\theta_\text{old}}$, PPO \citep{ppo} constrains the policy update within a proximal region of the old policy through the clipping mechanism. Specifically, PPO employs the following objective for policy optimization (we omit the KL regularization term hereinafter for brevity, as it is not the focus of this paper):

\begin{align}
\mathcal{J}_\text{PPO}(\theta) = \mathbb{E}_{ x \sim \mathcal{D},\, y \sim \pi_{\theta_\text{old}}( \cdot | x) }
\left[ \frac{1}{|y|} \sum_{t=1}^{|y|} 
\min \left( w_{t}(\theta) \widehat{A}_{t},  \, \mathrm{clip} \left( w_{t}(\theta), 1 - {\varepsilon}, 1 + {\varepsilon}\right) \widehat{A}_{t} \right)
\right],
\end{align}

The token importance ratio $w_{t}(\theta)$ is given by $
w_{t}(\theta) = \frac{ \pi_{\theta} (y_{t} | x, y_{<t}) }{ \pi_{\theta_\text{old}} (y_{t} | x,y_{<t})} $. Here, the advantage $\widehat{A}_{t}$ for $y_t$ is approximated using a distinct value model, and the parameter $\varepsilon$ determines the clipping threshold applied to the importance ratios.

A primary obstacle encountered when employing PPO arises from its substantial dependence on the value model. The value model commonly matches the policy model in scale, which results in significant resource demands for both memory and computation. Additionally, the success of the approach is intimately linked to how accurate the value predictions are. Developing a dependable value model is fundamentally difficult, and achieving robustness for tasks involving lengthy outputs and increased complexity escalates these difficulties further.

\paragraph{Group Relative Policy Optimization (GRPO)} GRPO \citep{grpo} bypasses the need for the value model by computing the relative advantage of each response within a group of responses to the same query. Specifically, GRPO optimizes the following objective:

\begin{align}
\label{equ:grpo}
\mathcal{J}_\text{GRPO}(\theta) = \mathbb{E}_{ x \sim \mathcal{D},\, \{y_i\}_{i=1}^G \sim \pi_{\theta_\text{old}}( \cdot | x) }
\left[ \frac{1}{G} \sum_{i=1}^{G} \frac{1}{|y_i|} \sum_{t=1}^{|y_i|} 
\min \left( w_{i,t}(\theta) \widehat{A}_{i,t},  \, \mathrm{clip} \left( w_{i,t}(\theta), 1 - {\varepsilon}, 1 + {\varepsilon}\right) \widehat{A}_{i,t} \right)
\right],
\end{align}

where $G$ is the number of generated responses to each query $x$ (i.e., the group size), and the importance ratio $w_{i,t}(\theta)$ and advantage $\widehat{A}_{i,t}$ of token $y_{i,t}$ are:

\begin{align}
    w_{i,t}(\theta)=\frac{ \pi_{\theta} (y_{i,t} | x, y_{i,<t}) }{ \pi_{\theta_\text{old}} (y_{i,t} | x,y_{i,<t})},\quad\ 
    \widehat{A}_{i,t} = \widehat{A}_{i} = \frac{r(x, y_i) - \mathrm{mean} \left( \{ r(x, y_i) \}_{i=1}^G \right) }{ \mathrm{std} \left( \{ r(x, y_i) \}_{i=1}^G \right) },
\end{align}

respectively, where all the tokens in $y_i$ share the same advantage as $\widehat{A}_{i}$.

\section{Motivation}
\label{sec:motivation}

As model parameters, sparsity (such as that found in Mixture-of-Experts architectures), and output lengths increase, reinforcement learning workflows require correspondingly larger rollout batch sizes to efficiently leverage hardware resources. To enhance the efficiency of sample usage, practitioners commonly subdivide these large rollout batches into several smaller mini-batches for the purpose of gradient computation. This approach unavoidably results in an off-policy learning context, since the generated responses $y$ are drawn from an earlier policy version $\pi_{\theta_\text{old}}$ and not from the most recently updated policy $\pi_\theta$. This scenario underscores the role of clipping mechanisms in methods like PPO and GRPO, which are introduced to limit the influence of samples that diverge excessively from the current policy during gradient calculations.

While mechanisms like clipping aim to manage this off-policy discrepancy, we identify a more fundamental issue in GRPO: \textit{its objective is ill-posed}. This problem becomes particularly acute when training large models on long-response tasks, leading to catastrophic model collapse. The ill-posed nature of the GRPO objective stems from a misapplication of importance sampling weights. The principle of importance sampling is to estimate the expectation of a function $f$ under a target distribution $\pi_\text{tar}$ by re-weighting samples drawn from a behavior distribution $\pi_\text{beh}$:

\begin{align}
\mathbb{E}_{z \sim \pi_\text{tar}} \left[ f(z) \right] = \mathbb{E}_{z \sim \pi_\text{beh}} \left[ \frac{ \pi_\text{tar} (z) }{ \pi_\text{beh} (z)} \, f(z) \right].
\end{align}

A key aspect of this approach is that it necessitates averaging across a large number of samples ($N \gg 1$) drawn from the behavior policy $\pi_\text{beh}$ to ensure that the importance weight $\frac{ \pi_\text{tar} (z) }{ \pi_\text{beh} (z)}$ properly adjusts for the divergence between target and behavior distributions.

By contrast, GRPO utilizes the importance weight $\frac{ \pi_{\theta} (y_{i,t} | x, y_{i,<t}) }{ \pi_{\theta_\text{old}} (y_{i,t} | x,y_{i,<t})}$ at every token position $t$. Because this weighting is applied to just a single realization $y_{i,t}$ from each next-token distribution $\pi_{\theta_\text{old}}( \cdot | x, y_{i,<t})$, it cannot adequately mitigate distribution mismatch as intended. Rather, it injects substantial variance into gradient estimates during training, which then aggregates across lengthy sequences and is further intensified by the use of clipping. Our empirical investigations indicate that such instability can drive the model toward catastrophic failure—an effect that is not readily reversible. Even after restoring earlier model checkpoints and conducting extensive hyperparameter searches (for instance, modifying clipping thresholds), lengthening output sequences, or changing reinforcement learning queries, recovery proves futile once collapse sets in.

This empirical evidence exposes a deeper problem at the core of GRPO’s methodology. Specifically, it highlights a crucial mismatch: \textbf{the granularity of the optimization objective needs to coincide with that of the reward}. Given that rewards are dispensed at the entire sequence level, attempts to impose off-policy adjustments on individual tokens are fundamentally flawed. Consequently, this observation prompts us to abandon token-level objectives in favor of leveraging importance weighting and optimization strategies that operate at the \textit{sequence level}.

\section{Algorithm}

\subsection{GSPO: Group Sequence Policy Optimization}

Although the token-level importance weight $\frac{ \pi_{\theta} (y_{i,t} | x, y_{i,<t}) }{ \pi_{\theta_\text{old}} (y_{i,t} | x,y_{i,<t})}$ poses challenges within GRPO, our findings indicate that, within language generation, the \textit{sequence-level} importance weight $\frac{ \pi_{\theta} (y | x) }{ \pi_{\theta_\text{old}} (y | x)}$ possesses a well-defined theoretical interpretation. Specifically, it quantifies the extent to which a response $y$ drawn from $\pi_{\theta_\text{old}} (\cdot | x)$ differs from that under $\pi_{\theta} (\cdot | x)$. This interpretation resonates naturally with the sequence-level reward structure and can be employed as a relevant gauge for the operation of the clipping mechanism.

Based on this straightforward observation, we propose the \textbf{Group Sequence Policy Optimization (GSPO)} algorithm. GSPO employs the following sequence-level optimization objective:

\begin{align}
\mathcal{J}_\text{GSPO} (\theta) =
\mathbb{E}_{ x \sim \mathcal{D},\, \{y_i\}_{i=1}^G \sim \pi_{\theta_\text{old}}( \cdot | x) }
\left[ 
\frac{1}{G} \sum_{i=1}^{G}
\min \left( s_{i}(\theta)  \widehat{A}_{i},  \, \mathrm{clip} \left( s_{i}(\theta), 1 - {\varepsilon}, 1 + {\varepsilon} \right) \widehat{A}_{i} \right) 
\right],
\label{equ:gspo}
\end{align}

where we adopt the group-based advantage estimation:

\begin{align}
\widehat{A}_{i} = \frac{r(x, y_i) - \mathrm{mean} \left( \{ r(x, y_i) \}_{i=1}^G \right) }{ \mathrm{std} \left( \{ r(x, y_i) \}_{i=1}^G \right) },
\end{align}

and define the importance ratio $s_{i}(\theta)$ based on sequence likelihood \citep{click}:

\begin{align}
s_{i}(\theta) = \left( \frac{ \pi_{\theta} (y_i | x) }{ \pi_{\theta_\text{old}} (y_i | x)} \right)^{\frac{1}{|y_i|}}
=
\exp \left( \frac{1}{|y_i|} \sum_{t=1}^{|y_i|} \log \frac{ \pi_{\theta} (y_{i,t} | x, y_{i,<t}) }{ \pi_{\theta_\text{old}} (y_{i,t} | x,y_{i,<t})} \right).
\end{align}

Consequently, GSPO implements clipping across whole responses rather than on a per-token basis, effectively preventing highly ``off-policy'' instances from contributing to gradient computation and aligning with both sequence-level reward assignment and optimization procedures. It should be noted that $s_{i}(\theta)$ is normalized for length, which serves to mitigate variance and maintain $s_{i}(\theta)$ within a consistent numerical scope. Without such normalization, changes in token likelihood for a small subset can cause substantial volatility in the sequence-level importance ratio, necessitating different clipping thresholds for responses of varying lengths. Additionally, it is worth mentioning that the clipping thresholds used in GSPO differ by orders of magnitude from those in earlier methods (such as GRPO) owing to fundamentally different formulations of importance ratios.

\subsection{Gradient Analysis}
\label{subsec:gradient}

We can derive the gradient of the GSPO objective as follows (clipping is omitted for brevity):

\begin{align}
\nabla_{\theta} \mathcal{J}_\text{GSPO} (\theta)
=&\ 
\nabla_{\theta} \mathbb{E}_{ x \sim \mathcal{D},\, \{y_i\}_{i=1}^G \sim \pi_{\theta_\text{old}}( \cdot | x) }
\left[ 
\frac{1}{G} \sum_{i=1}^{G} s_{i}(\theta) \widehat{A}_{i}
\right] \\
= &\ 
\mathbb{E}_{ x \sim \mathcal{D},\, \{y_i\}_{i=1}^G \sim \pi_{\theta_\text{old}}( \cdot | x) }
\left[ 
\frac{1}{G} \sum_{i=1}^{G}
s_{i}(\theta) \widehat{A}_{i}
\cdot \nabla_{\theta} \log s_{i}(\theta)
\right] \\
=&\ 
\mathbb{E}_{ x \sim \mathcal{D},\, \{y_i\}_{i=1}^G \sim \pi_{\theta_\text{old}}( \cdot | x) }
\left[ 
\frac{1}{G} \sum_{i=1}^{G} 
\left( \frac{ \pi_{\theta} (y_i | x) }{ \pi_{\theta_\text{old}} (y_i | x)} \right)^{\frac{1}{|y_i|}} \widehat{A}_{i} 
\cdot \frac{1}{|y_i|} \sum_{t=1}^{|y_i|} \nabla_{\theta} \log \pi_{\theta} (y_{i,t} | x, y_{i,<t}) 
\right].
\label{equ:gspo_grad}
\end{align}

For comparison, the gradient of the GRPO objective is as follows (note that $\widehat{A}_{i,t} = \widehat{A}_{i}$):

\begin{align}
\nabla_{\theta} \mathcal{J}_\text{GRPO} (\theta) 
=&\ 
\nabla_{\theta} \mathbb{E}_{ x \sim \mathcal{D},\, \{y_i\}_{i=1}^G \sim \pi_{\theta_\text{old}}( \cdot | x) }
\left[ \frac{1}{G} \sum_{i=1}^{G} \frac{1}{|y_i|} \sum_{t=1}^{|y_i|} 
w_{i,t}(\theta) \widehat{A}_{i,t}
\right] \\
=&\
\mathbb{E}_{ x \sim \mathcal{D},\, \{y_i\}_{i=1}^G \sim \pi_{\theta_\text{old}}( \cdot | x) }
\left[ \frac{1}{G} \sum_{i=1}^{G} \widehat{A}_{i} 
\cdot \frac{1}{|y_i|} \sum_{t=1}^{|y_i|} 
\frac{ \pi_{\theta} (y_{i,t} | x, y_{i,<t}) }{ \pi_{\theta_\text{old}} (y_{i,t} | x,y_{i,<t})} 
\nabla_{\theta} \log \pi_{\theta} (y_{i,t} | x, y_{i,<t})  
\right].
\end{align}

Consequently, the main difference between GSPO and GRPO concerns the \textit{manner in which gradients of token log likelihoods are weighted}. GRPO assigns each token a unique ``importance weight'' given by $\frac{ \pi_{\theta} (y_{i,t} | x, y_{i,<t}) }{ \pi_{\theta_\text{old}} (y_{i,t} | x,y_{i,<t})}$, resulting in non-uniform weighting across tokens. These weights, which can range within $(0, 1+\varepsilon]$ (when $\widehat{A}_{i} >0$) or $[1-\varepsilon, +\infty)$ (when $\widehat{A}_{i} < 0$), may substantially affect the training process as they are non-trivial and their influence can become pronounced over time, potentially causing unpredictable behaviors. On the other hand, GSPO applies uniform weighting to all tokens in a given response, thereby removing the source of instability inherent in GRPO.

\subsection{GSPO-token: A Token-level Objective Variant}

In scenarios like multi-turn RL, we may desire a finer-grained advantage adjustment than the sequence level. To this end, we introduce a token-level objective variant of GSPO, namely \textbf{GSPO-token}, to allow token-wise advantage customization:

\begin{align}
\mathcal{J}_\text{GSPO-token}(\theta)
=
\mathbb{E}_{ x \sim \mathcal{D},\, \{y_i\}_{i=1}^G \sim \pi_{\theta_\text{old}}( \cdot | x) }
\left[ 
\frac{1}{G} \sum_{i=1}^{G} \frac{1}{|y_i|} \sum_{t=1}^{|y_i|} 
\min \left( s_{i,t}(\theta) \widehat{A}_{i,t},  \, \mathrm{clip} \left( s_{i,t}(\theta), 1 - {\varepsilon}, 1 + {\varepsilon} \right) \widehat{A}_{i,t} \right) 
\right],
\label{equ:gspo-token}
\end{align}

where

\begin{align}
s_{i,t}(\theta) 
= \mathrm{sg} \left[ s_{i}(\theta) \right]  \cdot \frac{ \pi_{\theta} (y_{i,t} | x, y_{i,<t}) }{ \mathrm{sg} \left[ \pi_{\theta} (y_{i,t} | x, y_{i,<t}) \right] },
\end{align}

and $\mathrm{sg}[\cdot]$ denotes only taking the numerical value but stopping the gradient, corresponding to the \texttt{detach} operation in PyTorch. The gradient of GSPO-token can be derived as:

\begin{align}
\nabla_{\theta} \mathcal{J}_\text{GSPO-token}(\theta)
=&\ 
\nabla_{\theta} \mathbb{E}_{ x \sim \mathcal{D},\, \{y_i\}_{i=1}^G \sim \pi_{\theta_\text{old}}( \cdot | x) }
\left[ 
\frac{1}{G} \sum_{i=1}^{G}
\frac{1}{|y_i|} \sum_{t=1}^{|y_i|} 
s_{i,t}(\theta) \widehat{A}_{i,t}
\right] \\
=&\ 
\mathbb{E}_{ x \sim \mathcal{D},\, \{y_i\}_{i=1}^G \sim \pi_{\theta_\text{old}}( \cdot | x) }
\left[ 
\frac{1}{G} \sum_{i=1}^{G} s_i (\theta) \cdot \frac{1}{|y_i|} \sum_{t=1}^{|y_i|} 
\widehat{A}_{i,t} \frac{ \nabla_{\theta} \pi_{\theta} (y_{i,t} | x, y_{i,<t}) }{ \pi_{\theta} (y_{i,t} | x, y_{i,<t}) }
\right] \\
=&\ 
\mathbb{E}_{ x \sim \mathcal{D},\, \{y_i\}_{i=1}^G \sim \pi_{\theta_\text{old}}( \cdot | x) }
\left[ 
\frac{1}{G} \sum_{i=1}^{G}  \left( \frac{ \pi_{\theta} (y_i | x) }{ \pi_{\theta_\text{old}} (y_i | x)} \right)^{\frac{1}{|y_i|}}
\cdot \frac{1}{|y_i|} \sum_{t=1}^{|y_i|}
\widehat{A}_{i,t} \nabla_{\theta} \log \pi_{\theta} (y_{i,t} | x, y_{i,<t})  
\right].
\label{equ:gspo-token_grad}
\end{align}

It is important to observe that the expression $\frac{ \pi_{\theta} (y_{i,t} | x, y_{i,<t}) }{ \mathrm{sg} \left[ \pi_{\theta} (y_{i,t} | x, y_{i,<t}) \right] }$ evaluates to 1, meaning that $s_{i,t}(\theta)$ yields the same numerical result as $s_{i}(\theta)$. If we examine Equation~\eqref{equ:gspo} and~\eqref{equ:gspo-token}, along with Equation~\eqref{equ:gspo_grad} and~\eqref{equ:gspo-token_grad}, it becomes evident that GSPO-token and GSPO correspond to identical quantities regarding the optimization objective, clipping condition, and theoretical gradient—provided that all tokens within the response $y_i$ are assigned the same advantage (i.e., $\widehat{A}_{i,t} = \widehat{A}_{i}$). Nevertheless, GSPO-token offers greater adaptability, enabling the adjustment of advantages at the token level.

\section{Experiments and Discussion}

\subsection{Empirical Results}

We utilize a cold-start configuration based on Qwen3-30B-A3B-Base for fine-tuning, presenting both the reward trajectories during training and the model’s progression on three benchmarks: AIME'24 (reporting average Pass@1 across 32 samples), LiveCodeBench (202410-202502, with average Pass@1 over 8 samples), and CodeForces (using Elo Rating). The RL training regime involves splitting each batch of rollout data into four mini-batches, facilitating multiple gradient update steps per batch. For GSPO, the left and right clipping bounds in Equation~\eqref{equ:gspo} are selected as 3e-4 and 4e-4, respectively. Our experiments include GRPO as a baseline, with the left and right clipping intervals in Equation~\eqref{equ:grpo} set to 0.2 and 0.27, ensuring these values are well-tuned for a just comparison. Importantly, successful convergence of MoE RL with GRPO depends on the Routing Replay approach, which is further elaborated in \S~\ref{subsec:routing_replay}; in contrast, \textit{GSPO eliminates the necessity for this strategy}.

As illustrated in Figure~\ref{fig:results}, the training process utilizing GSPO remains consistently stable. Our findings indicate that \textbf{GSPO facilitates ongoing enhancements in performance by increasing the training compute, periodically refreshing the query set, and lengthening the generation outputs}. In addition, GSPO exhibits greater training efficiency compared to GRPO, obtaining higher accuracy and improved benchmark outcomes given identical amounts of training compute and query usage. Lastly, the successful implementation of GSPO in RL training for the latest Qwen3 models offers compelling evidence of GSPO's effectiveness in harnessing the benefits of RL scaling in large language model development.

\subsection{Curious Observation on Clipping Fractions}

One fundamental difference between GSPO and GRPO lies in GSPO's method of clipping at the response level rather than at the token level. As illustrated in Figure~\ref{fig:clipping}, there exists a disparity of two orders of magnitude regarding the proportion of clipped tokens between GSPO and GRPO, and this gap remains consistent even when the clipping ranges are modified. Interestingly, GSPO clips far more tokens—thereby reducing the number utilized for model training or gradient estimation—but nonetheless demonstrates greater training efficiency compared to GRPO. This surprising outcome, where extensive token clipping correlates with improved training efficiency, underscores that GRPO's token-wise gradient estimates are fundamentally noisier and less effective in exploiting samples. GSPO's reliance on sequence-level clipping, therefore, yields a more dependable and impactful learning signal.

\subsection{Benefit of GSPO for MoE Training}
\label{subsec:routing_replay}

\paragraph{Background}
In contrast to RL training with dense architectures, the inherently sparse activation pattern characteristic of MoE models presents distinct challenges regarding training stability. Specifically, our observations indicate that when utilizing the GRPO algorithm, the MoE model’s \textit{expert-activation volatility} can disrupt RL training convergence. Concretely, the set of experts chosen to generate the same response may undergo substantial shifts following each gradient update. As an illustration, for the 48-layer Qwen3-30B-A3B-Base architecture, around 10\% of the experts selected by the updated policy $\pi_{\theta}$ for a given rollout differ from those selected by the previous policy $\pi_{\theta_\text{old}}$ after a single RL gradient step.

This effect intensifies as the depth of MoE models increases, which causes the token-level importance ratios, defined as $w_{i,t}(\theta) = \frac{ \pi_{\theta} (y_{i,t} | x, y_{i,<t}) }{ \pi_{\theta_\text{old}} (y_{i,t} | x,y_{i,<t})}$, to become highly unstable—often invalidating their utility, as elaborated in \S~\ref{sec:motivation} and~\ref{subsec:gradient}. As a result, this instability significantly impedes the standard convergence behavior expected during RL training.

\paragraph{Our Previous Approach} In our earlier work, we utilized the \textbf{Routing Replay} training method to address this issue. This involved recording the set of experts activated by $\pi_{\theta_\text{old}}$ and then “replaying” these expert selections during computation of the importance ratios $w_{i,t}(\theta) = \frac{ \pi_{\theta} (y_{i,t} | x, y_{i,<t}) }{ \pi_{\theta_\text{old}} (y_{i,t} | x, y_{i,<t})}$ under $\pi_{\theta}$. Consequently, when evaluating each token $y_{i,t}$, both $\pi_{\theta} (y_{i,t} | x, y_{i,<t})$ and $\pi_{\theta_\text{old}} (y_{i,t} | x, y_{i,<t})$ are routed through an identical set of network activations, which allows us to recover stability in token-level importance ratios and facilitate optimization within a consistently activated network throughout gradient updates. As illustrated in Figure~\ref{fig:routing_replay}, Routing Replay is critical to achieving stable convergence during GRPO training for MoE models.

\paragraph{Benefit of GSPO} While Routing Replay facilitates the convergence of GRPO training for MoE models, it does so at the expense of increased memory usage and communication, and it may also constrain the MoE model's usable capacity. GSPO, as illustrated in Figure~\ref{fig:results}, sidesteps the reliance on Routing Replay by conventionally calculating importance ratios $s_i(\theta)$, achieving steady optimization and standard convergence. The central realization is that GSPO operates exclusively with the sequence likelihood (specifically, $\pi_{\theta} (y_i | x)$) rather than being affected by token-level likelihoods ($\pi_{\theta} (y_{i,t} | x, y_{i,<t})$). Since MoE models inherently preserve their language modeling proficiency, the sequence likelihood remains relatively stable. Therefore, GSPO offers a fundamental solution to the instability of expert activation in MoE models, eliminating the necessity for elaborate methods such as Routing Replay. This contributes both to more robust and straightforward training procedures and to allowing the model to utilize its full capacity unhindered.

\subsection{Benefit of GSPO for RL Infrastructure}

Due to precision differences present between training systems such as Megatron and inference frameworks like SGLang or vLLM, it is common practice to recompute the likelihoods of sampled outputs according to the old policy $\pi_{\theta_\text{old}}$ using the original training engine. In contrast, GSPO leverages only sequence-level likelihoods in its optimization process, rather than relying on token-level measurements, which inherently grants greater robustness against such precision variances. As a result, GSPO permits direct utilization of likelihood values generated by the inference engine during optimization, effectively eliminating the requirement for likelihood recomputation via the training engine. This capability is particularly advantageous in contexts such as partial rollout, multi-turn reinforcement learning, and architectures where training and inference are handled by separate systems.

\section{Conclusion}

Group Sequence Policy Optimization (GSPO) is introduced as a novel reinforcement learning technique specifically for large language model training. Building upon the importance sampling framework, GSPO utilizes sequence likelihood to determine importance ratios, implementing sequence-level procedures for clipping, reward assignment, and optimization. In comparison to GRPO, GSPO achieves markedly enhanced stability during training, greater efficiency, and improved overall performance, showing particular advantages when applied to large-scale RL scenarios involving MoE models. This contribution has been pivotal in enabling significant advancements in the recent Qwen3 models. Leveraging GSPO's scalability, future work aims to further expand the capabilities of RL, anticipating substantial progress in the development of intelligent systems.

\end{document}